{\scriptsize
    \begin{minipage}[t]{0.5\textwidth}
        \begin{tabularx}{\linewidth}{|p{0.3\textwidth}|X|}
        \hline 
        Notions et contenus & Capacités exigibles \\
        \hline 
        \multicolumn{2}{|l|}{1.4. Exemple de dispositif interférentiel par division d'amplitude : }  \\
        \multicolumn{2}{|l|}{ interféromètre de Michelson} \\
        \hline 
        Interféromètre de Michelson équivalent à une lame d'air éclairée par une source spatialement étendue. 
         & Justifier les conditions d'observation des franges d'égale inclinaison, le lieu de localisation des franges étant admis. Établir et utiliser l'expression de l'ordre d'interférences en fonction de l'épaisseur de la lame, l'angle d'incidence et la longueur d'onde. \\
        \hline 
        Localisation des franges. Franges d'égale inclinaison. &
        Décrire et mettre en œuvre les conditions d'éclairage et d'observation adaptées à l'utilisation d'un interféromètre de Michelson en lame d'air. Mesurer l'écart en longueur d'onde d'un doublet et la longueur de cohérence d'une radiation. Interpréter des observations faites en lumière blanche avec l'interféromètre de Michelson en configuration lame d'air.  \\
        \hline
        \end{tabularx}
\end{minipage}%
\begin{minipage}[t]{0.5\textwidth}
    \begin{tabularx}{\linewidth}{|X|X|}
    \hline 
    Interféromètre de Michelson équivalent à un coin d'air éclairé par une source spatialement étendue. 
   & Justifier les conditions d'observation des franges d'égale épaisseur, le lieu de localisation des franges étant admis. Utiliser l'expression donnée de la différence de marche en fonction de l'épaisseur pour exprimer l'ordre d'interférences. \\
   \hline
   Localisation des franges. Franges d'égale épaisseur.   & Décrire et mettre en œuvre les conditions d'éclairage et d'observation adaptées à l'utilisation d'un interféromètre de Michelson en coin d'air. Caractériser la géométrie d'un objet ou l'indice d'un milieu à l'aide d'un interféromètre de Michelson. Interpréter des observations faites en lumière blanche avec l'interféromètre. \\
    \hline 
    \end{tabularx}
\end{minipage}
}
